\documentclass[11pt]{article}

\usepackage[a4paper,margin=1in]{geometry}
\usepackage{amsmath,amssymb,amsthm}
\usepackage{booktabs}
\usepackage{longtable}
\usepackage{microtype}
\usepackage{url}

\newtheorem{theorem}{Theorem}
\newtheorem{lemma}[theorem]{Lemma}
\newcommand{\dd}{\,\mathrm d}

\title{Exact Four-Root Interval Sum-of-Squares Certificates\\
for the Remaining Six-Vertex Cases}
\author{}
\date{August 25, 2026}

\begin{document}
\maketitle

\begin{abstract}
We record direct arbitrary-graphon proofs of the conjectured chromatic lower
bound for sixteen previously open or partial Atlas graphs.  The new input is
an exact rational interval sum-of-squares identity on four labelled vertices
and one integrated branch vertex.  Every identity is independently rebuilt
from the Atlas graph and checked coefficient by coefficient.  Positivity is
certified by rational Gram factors and strictly diagonally dominant rational
corrections.  The argument does not assume that an optimizer has a special
form and does not use flag-algebra inference rules.  Together with the
separate exact house certificate and the normalized-link house-cone theorem,
this closes eighteen of the twenty cases that were unresolved at the start of
this research pass.  Atlas 130 and Atlas 203 retain only proper-interval
results and therefore remain unclassified.
\end{abstract}

\section{Statement}

For a graph $H$ on $h$ vertices with chromatic polynomial $P_H$ put
\[
 \Phi_H(p)=(1-p)^hP_H\!\left(\frac1{1-p}\right).
\]
The required density interval is $p\in[1-1/(\chi(H)-1),1]$.

\begin{theorem}\label{thm:classification}
For every graphon $W$, the inequality
\[
             t(H,W)\geq \Phi_H(t(K_2,W))
\]
holds throughout the required density interval for each
\[
\begin{split}
H\in\{&118,122,124,147,151,153,157,168,169,171,174,\\
       &181,185,188,194,199\}.
\end{split}
\]
Here a number denotes the corresponding NetworkX Graph Atlas graph.
\end{theorem}

Atlas 43 is proved separately by the exact rooted identity in
\path{notes/atlas43_exact_rooted_sos.tex}.  Atlas 196 is the cone
$K_1\vee H_{43}$ and hence follows unconditionally from Atlas 43 by the
normalized-link theorem in \path{notes/open_conditional_implications.tex}.
Thus those two graphs are also positive classifications.

\section{The direct rooted identity}

Fix four labelled variables $x=(x_0,x_1,x_2,x_3)$ and one branch variable
$y$.  For each graph $F$ having those four labels and the one branch, let
\[
 f_F(x)=\int_\Omega\prod_{ij\in E(F)}W(z_i,z_j)\dd y,
 \qquad z_i=x_i\ (0\leq i<4),\quad z_4=y.
\]
The certificates use every branch-to-label edge pattern and labelled graphs
with at most two label-to-label edges.  Put these evaluations into a vector
$f(x)$.  Permuting the four labels gives an $S_4$ action.  Integer Young
symmetrizers split $f$ into the five rational slices
\[
                  \lambda\in\{4,31,22,211,1111\}.
\]
Write the corresponding evaluation vectors as $v_\lambda(x)$.

For an interval $[a,b]$, set
\[
                  p=a+(b-a)s,\qquad 0\leq s\leq1,
\]
and $\widehat v_\lambda=(v_\lambda,sv_\lambda)$.  Each JSON certificate
gives rational positive-semidefinite matrices $Q_{0,\lambda}$ and
$Q_{1,\lambda}$ for which exact expansion gives
\begin{equation}\label{eq:master}
\boxed{
 t(H,W)-\Phi_H(p)=
 24\sum_\lambda\left(
  \int_{\Omega^4}\widehat v_\lambda^{\mathsf T}
     Q_{0,\lambda}\widehat v_\lambda\dd x
 +s(1-s)\int_{\Omega^4}v_\lambda^{\mathsf T}
     Q_{1,\lambda}v_\lambda\dd x\right).}
\end{equation}
The factor 24 is the order of $S_4$.  This is simply a symmetry-reduced
way to write a sum of squares; it has no sign effect.

To see directly why \eqref{eq:master} is a graphon identity, expand a
quadratic term.  A product $f_F(x)f_G(x)$ uses two independent branch
variables and therefore is the homomorphism density of the graph obtained
by identifying equally labelled vertices of $F$ and $G$.  The result has at
most six vertices.  Any disjoint isolated-edge components are removed using
the exact fixed-density relation
\[
                 t(J\sqcup mK_2,W)=p^m t(J,W).
\]
After substituting $p=a+(b-a)s$, equality of every coefficient of every
unlabelled graph proves \eqref{eq:master} for arbitrary $W$.

The numerical SDP was used only to discover Gram matrices.  In a stored
certificate each final Gram block has the form
\[
             Q=D^{-2}LZL^{\mathsf T},\qquad Z=I+C,
\]
with integer $L,D$ and rational symmetric $C$.  The checker verifies in
exact fraction arithmetic that $Z=I+C$ has positive diagonal and is strictly
diagonally dominant.  It is therefore positive definite, so every $Q$ in
\eqref{eq:master} is positive semidefinite.  Exact complements of kernels
seen at balanced Turan equality graphons were used to reduce the matrices.
This is only a rational change to a smaller Gram-coordinate space: it is
not a restriction on $W$.

\begin{lemma}\label{lem:interval}
An exact certificate satisfying the coefficient and diagonal-dominance
checks proves $t(H,W)\geq\Phi_H(p)$ for every graphon $W$ and every
$p\in[a,b]$.
\end{lemma}

\begin{proof}
All integrals on the right of \eqref{eq:master} are nonnegative quadratic
forms.  Also $s(1-s)\geq0$ on $[0,1]$.
\end{proof}

\section{Certificates and interval coverage}

The exact files used in Theorem~\ref{thm:classification} are listed below.
The last column is a decimal lower bound for the smallest exact strict
diagonal-dominance margin among all corrected Gram blocks.  It is included
only as a readability aid; the JSON stores the exact rational value.

\begin{longtable}{@{}rllr@{}}
\toprule
Atlas & interval & certificate basename & minimum margin $>$\\
\midrule
\endhead
118 & $[1/2,1]$ & \path{atlas118_exact_interval_sos.json} & $0.9751$\\
122 & $[1/2,1]$ & \path{atlas122_exact_interval_sos.json} & $0.9764$\\
124 & $[1/2,1]$ & \path{atlas124_exact_interval_sos.json} & $0.9805$\\
147 & $[1/2,2/3]$ & \path{atlas147_exact_lower_interval_sos.json} & $0.9665$\\
147 & $[2/3,1]$ & \path{atlas147_exact_upper_interval_sos.json} & $0.9948$\\
151 & $[1/2,2/3]$ & \path{atlas151_exact_lower_interval_sos.json} & $0.9929$\\
151 & $[2/3,1]$ & \path{atlas151_exact_upper_interval_sos.json} & $0.9826$\\
153 & $[37/60,2/3]$ & \path{atlas153_exact_middle_interval_sos.json} & $0.9961$\\
153 & $[2/3,1]$ & \path{atlas153_exact_upper_interval_sos.json} & $0.9970$\\
157 & $[2/3,1]$ & \path{atlas157_exact_interval_sos.json} & $0.9789$\\
168 & $[1/2,2/3]$ & \path{atlas168_exact_lower_interval_sos.json} & $0.9922$\\
168 & $[2/3,1]$ & \path{atlas168_exact_upper_interval_sos.json} & $0.9943$\\
169 & $[2/3,1]$ & \path{atlas169_exact_interval_sos.json} & $0.9746$\\
171 & $[5/8,2/3]$ & \path{atlas171_exact_middle_interval_sos.json} & $0.9911$\\
171 & $[2/3,1]$ & \path{atlas171_exact_upper_interval_sos.json} & $0.9923$\\
174 & $[7/12,2/3]$ & \path{atlas174_exact_middle_interval_sos.json} & $0.8831$\\
174 & $[2/3,1]$ & \path{atlas174_exact_upper_interval_sos.json} & $0.9833$\\
181 & $[2/3,1]$ & \path{atlas181_exact_interval_sos.json} & $0.9726$\\
185 & $[2/3,1]$ & \path{atlas185_exact_interval_sos.json} & $0.9930$\\
188 & $[3/5,2/3]$ & \path{atlas188_exact_middle_interval_sos.json} & $0.9761$\\
188 & $[2/3,1]$ & \path{atlas188_exact_upper_interval_sos.json} & $0.9995$\\
194 & $[2/3,1]$ & \path{atlas194_exact_interval_sos.json} & $0.9935$\\
199 & $[2/3,1]$ & \path{atlas199_exact_interval_sos.json} & $0.9901$\\
\bottomrule
\end{longtable}

All basenames in the table lie in \path{experiments/}.  For Atlas 147,
151, and 168, the two closed intervals meet and cover $[1/2,1]$.  The
single certificates for 118, 122, and 124 cover $[1/2,1]$ directly.  The
single certificates for 157, 169, 181, 185, 194, and 199 cover their whole
required interval $[2/3,1]$ directly.

Four low-density pieces come from the previously accepted arbitrary-graphon
self-amalgam theorem in
\path{notes/four_self_amalgam_partial_ranges.tex}:
\[
\begin{array}{c|c}
H&\hbox{previously proved range}\\ \hline
153 &[1/2,\,0.6181289270\ldots]\\
171 &[1/2,\,0.6309517719\ldots]\\
174 &[1/2,\,0.5934675772\ldots]\\
188 &[1/2,\,0.6139911259\ldots].
\end{array}
\]
These overlap the exact middle intervals because
\begin{align*}
 37/60&<0.6181289270\ldots,&
 5/8&<0.6309517719\ldots,\\
 7/12&<0.5934675772\ldots,&
 3/5&<0.6139911259\ldots.
\end{align*}
Their unions with the middle and upper certificates therefore cover
$[1/2,1]$.  Lemma~\ref{lem:interval} proves every exact piece, and the
stated interval unions prove Theorem~\ref{thm:classification}.

\section{Exact audit}

The independent checker is
\path{experiments/s4_interval_verify.py}.  Given only one JSON file, it
rebuilds the Atlas graph and chromatic target, reconstructs the rooted
products and the five integer Young slices, checks the basis and dimensions,
checks all strict diagonal-dominance inequalities over $\mathbb Q$, and
then verifies all 407 coefficient equations with FLINT integer and rational
matrices.  Thus the numerical solver and the floating certificate used in
discovery are outside the trusted proof path.

For example, the audit command has the form
\begin{verbatim}
.venv/Scripts/python.exe experiments/s4_interval_verify.py \
  experiments/atlas181_exact_interval_sos.json
\end{verbatim}
and reports 407 exact equations before exiting successfully.  The house
certificate has a different three-root basis and is audited independently by
\path{experiments/house_verify_rational.py}.

\subsection{Validation environment and commands}

The accepted audit was run with Python~3.9 and the following package
versions:
\begin{center}
\begin{tabular}{ll@{\qquad}ll}
\toprule
package & version & package & version\\
\midrule
\texttt{cvxpy} & \texttt{1.7.5} & \texttt{networkx} & \texttt{3.2.1}\\
\texttt{numpy} & \texttt{2.0.2} & \texttt{python-flint} & \texttt{0.6.0}\\
\texttt{scipy} & \texttt{1.13.1} & \texttt{sympy} & \texttt{1.14.0}\\
\bottomrule
\end{tabular}
\end{center}
These versions are recorded for reproducibility, not used as mathematical
assumptions.  From the \path{exhaustive_research/} directory, the two checker
entry points are
\begin{verbatim}
python experiments/house_verify_rational.py
python experiments/s4_interval_verify.py \
  experiments/atlas181_exact_interval_sos.json
\end{verbatim}
On PowerShell, all committed four-root certificates can be checked by
\begin{verbatim}
Get-ChildItem experiments -Filter 'atlas*_exact_*interval_sos.json' |
  ForEach-Object {
    python experiments/s4_interval_verify.py $_.FullName
    if ($LASTEXITCODE -ne 0) {
      throw "certificate audit failed: $($_.Name)"
    }
  }
\end{verbatim}
The committed JSON files are the exact proof certificates.  Numerical
\texttt{.npz} workspaces were used only to discover the Gram matrices; they
are neither trusted by these commands nor required to validate the proof.

\section{The two cases not closed}

The same exact method gives two useful but incomplete arbitrary-graphon
results:
\begin{center}
\begin{tabular}{ccl}
\toprule
$H$ & proved interval & certificate\\
\midrule
130 & $[2/3,1]$ & \path{experiments/atlas130_exact_upper_interval_sos.json}\\
203 & $[3/4,1]$ & \path{experiments/atlas203_exact_upper34_interval_sos.json}\\
\bottomrule
\end{tabular}
\end{center}
Atlas 130 still lacks $[1/2,2/3)$, and Atlas 203 still lacks
$[2/3,3/4)$.  These are proper-range theorems, not complete positive
classifications.  No exact counterexample was found for either graph.  The
final search job was a CVXOPT fixed-margin feasibility solve for Atlas 203 on
$[2/3,7/10]$ in the degree-two four-root cone.  It ran for approximately
42 minutes, ended with a solver failure, and wrote no candidate certificate;
this numerical failure has no mathematical sign implication.

\end{document}
